\section{范数规模、实际系数代价与低秩障碍}\label{sec:barriers}
\subsection{频率分离给出的上界}
固定常数 $B_\psi,B_\kappa<\infty$ 使
\[
0\le\psi_+(v)\le B_\psi(1+|v|)^{-2},\qquad
|\kappa(y)|\le B_\kappa(1+|y|)^{-2}.
\]
\begin{lemma}\label{lem:upper}
对 $Q\ge2$，
\[
\norm A\le4B_\psi(H+4Q^2)+3Q^2.
\]
\end{lemma}
\begin{proof}
不同约分分数 $a/q$ 在单位圆上至少相隔 $(4Q^2)^{-1}$，跨越零点的间隔亦然。定义在全部整数上的卷积核
$\kappa(h/H)\sum_{q,a}\e(ah/q)$ 的 Fourier 乘子为
\[
m(\theta)=H\sum_{q,a}\sum_{k\in\ZZ}
 \psi_+\big(H(a/q-\theta-k)\big)\ge0 .
\]
这由固定 Schwartz 核的 Poisson 求和得到。右边各点的实轴间距至少为
$s=H/(4Q^2)$；每个半开单位区间含点不超过 $1+1/s$。分割正负半轴并求和
$(1+j)^{-2}$，得到 $m(\theta)\le4B_\psi H(1+1/s)$。
压缩到 $I_x$ 后范数不增加。命题~\ref{prop:decomp}给出
$\norm{A+\diag\Delta}\le\norm T$，而
$\max\Delta<\sum_{Q<q\le2Q}q\le3Q^2$，故结论成立。
\end{proof}
由于 $H<Q^2$，实际尺度上 $\norm A\ll_\psi Q^2$。这一上界不使用 $\beta,w$ 的消去。

\subsection{所有实际行的能量}
\begin{proposition}\label{prop:energy}
对固定核和所有充分大的 $x$，一致于每个 $u\in I_x$，
\[
G(u)\asymp_\psi \frac{HQ^2}{\log Q}.
\]
\end{proposition}
\begin{proof}
连续性与 $\kappa(0)=1$ 给出固定 $0<\rho,b\le1$，使 $|\kappa(y)|\ge b$ 于 $|y|\le\rho$。在 $q\nmid u$ 时令
\[
S_q(u)=\sum_{\substack{t\in I_x\\t\ne u,\ q\nmid t}}|\kappa((u-t)/H)|^2,\quad
C_q(u)=\sum_{\substack{t\in I_x\\t\ne u,\ q\mid u-t}}|\kappa((u-t)/H)|^2.
\]
平方展开得到精确式
\[
G(u)=\sum_{\substack{q\in\QQ\\q\nmid u}}
 a_q^2\big(S_q(u)+(q-1)(q-3)C_q(u)\big).
\]
积分比较给出 $S_q(u)\le2B_\kappa^2H/3$、
$C_q(u)\le2B_\kappa^2H/(3q)$。因为 $1+(q-1)(q-3)/q\le q$，
每个素数的上界为 $(2B_\kappa^2/3)a_q^2qH$。

当 $N-1\ge2\rho H$ 且 $H\ge2Q/\rho$ 时，每行至少一侧含
$\lfloor\rho H/q\rfloor\ge\rho H/(2q)$ 个非零的同余差。这些点仍为单位剩余类，
且 $(q-2)^2\ge q^2/9$，所以该素数的贡献至少为
$(\rho b^2/18)a_q^2qH$。

任何 $u\le x=Q^3$ 至多有两个大于 $Q$ 的不同素因子。因此从
$S_Q=\sum_{q\in\QQ}a_q^2q$ 中因 $q\mid u$ 删除的质量至多 $9Q$。
素数定理和分部求和给出
$S_Q\sim(3/2)Q^2/\log Q$。在一个共同阈值后 $S_Q\ge18Q$，从而
\[
\frac{\rho b^2}{36}HS_Q\le G(u)\le\frac{2B_\kappa^2}{3}HS_Q
\]
对全部实际行成立。
\end{proof}
区间上界 $u\le Q^3$ 在这个论证中不可删除；这不是任意巨大起点上的同长度区间定理，也不是任意坐标前缀的行能量定理。

\subsection{常向量见证与实际系数}
\begin{theorem}[系数无关范数方法的障碍]\label{thm:normbarrier}
对第~\ref{sec:operator}节的固定数据，
\[
\norm A\ge(3/2+o(1))\frac{Q^2}{\log Q},\qquad
\norm w^2\sim\frac x2\log x,\qquad
\norm\beta^2\ge\left(\frac12\log\frac{27}{22}+o(1)\right)\frac x{\log x}.
\]
因此
\[
\liminf_{x\to\infty}\frac{\log Q}{xQ^2}
 \norm A\,\norm\beta\,\norm w
\ge\frac34\sqrt{\log(27/22)}>0.
\]
任何仅用 $B_x\norm\beta\norm w$ 且 $B_x\ge\norm A$ 的上界证明，都不能单独支付固定正幂次节省。
\end{theorem}
\begin{proof}
令 $f=N^{-1/2}\1_{I_x}$。未删对角的正形式为
\[
\ip f{(A+\diag\Delta)f}
=\frac1N\sum_q q\int\psi_+(v)V_q(v/H)\,dv ,
\]
其中 $V_q(\alpha)$ 是各单位剩余类上
$\sum_{n\equiv r(q),\,n\in I_x}\e(\alpha n)$ 减去其单位剩余类均值后的平方和，
不再除以 $q-1$。写 $I_x=\{a+1,\ldots,a+N\}$、$N=mq+s$，$0\le s<q$，并定义
\[
G_m(y)=\sum_{k=0}^{m-1}\e(ky),\qquad
S_j(\alpha)=\e(\alpha(a+j))
\bigl(G_m(\alpha q)+\1_{j\le s}\e(\alpha mq)\bigr),\quad 1\le j\le q.
\]
$S_j$ 就是剩余类 $a+j$ 上的和，包含 $m$ 个完整周期加至多一个末项。当
$0<|\alpha|q\le1/2$ 时，
\[
|G_m(\alpha q)|\le(2|\alpha|q)^{-1},\qquad
|\e(\alpha j)-1|\le2\pi|\alpha|q,\quad 1\le j\le q.
\]
故各和距一个共同中心不超过 $\pi+1$，在 $\alpha=0$ 时不超过 $1$。
减去均值最小化平方距离，得到 $V_q(\alpha)\le q(\pi+1)^2$。
由于核的频率支撑固定且 $2Q/H<1/2$ 最终成立，整个积分都在这个范围内。
正形式于是为 $O(Q^3/(N\log Q))=O(1/\log Q)$。
另一方面，单位整数数目 $N_q=N(1-1/q)+O(1)$ 给出
\[
\ip f{Af}=-\sum_{q\in\QQ}(q-2)+O(1/\log Q).
\]
素数定理与分部求和推出所述算子下界。

为控制 $w$，令
$c(p)=((p-1)/(p-2))^2-1=O(1/p)$，并在无平方因子整数上乘法延拓。
展开非负 Euler 乘积，去掉至多为 $1$ 的掩码，并扩大区间，得到
\[
\norm{b_x^{(z)}}^2
\le xW_z^2\prod_{p>z}(1+c(p)/p)\ll xW_z^2.
\]
乘积收敛，因为 $c(p)/p=O(p^{-2})$。Mertens 乘积定理给出
$W_z\ll\log z\ll_K\log\log x$。素数定理又给出
$\sum_{u\in I_x}\Lambda(u+2)^2\sim(x/2)\log x$；
真素数幂贡献为低阶项。三角不等式的双向范数比较即证明 $w$ 的渐近式，
不需假定 $b_x^{(z)}$ 与 $\Lambda$ 正交。

为控制 $\beta$，取素数
$x^{2/5}<p\le x^{9/20}$ 和 $x/(2p)<r\le x/p$。
最终 $p<r$ 且二者均大于 $U$，故 $n=pr\in I_x$、$\Lambda(n)=0$，
且 $n$ 在 $U$ 以下唯一的约数是 $1$，所以 $\beta(n)=-1$。
较小素因子唯一，计数不重复。普通二倍区间的素数定理在此范围内一致，故这些乘积数目为
\[
(1+o(1))\frac x2\sum_{x^{2/5}<p\le x^{9/20}}\frac1{p\log(x/p)}
=\left(\frac12\log\frac{27}{22}+o(1)\right)\frac x{\log x}.
\]
最后一步由分部求和和
$\int_{2/5}^{9/20}[a(1-a)]^{-1}\,da=\log(27/22)$ 得到。
将三个下界相乘即得结论。
\end{proof}

\begin{corollary}\label{cor:normalized}
对所定义的行能量归一化，
\[
H^{-1}\ll_\psi\norm Z\ll_\psi(\log Q)/H.
\]
而 $\norm Z\,\norm{D_G^{1/2}\beta}\,\norm{D_G^{1/2}w}$
仍至少为 $c_\psi xQ^2/\log Q$，其中 $c_\psi>0$。
\end{corollary}
\begin{proof}
由 $A=D_G^{1/2}ZD_G^{1/2}$，
$\norm A/G_{\max}\le\norm Z\le\norm A/G_{\min}$。
两个加权向量各自至少引入 $\sqrt{G_{\min}}$，
故其范数乘积上界表达至少为
$(G_{\min}/G_{\max})\norm A\norm\beta\norm w$。
应用命题~\ref{prop:energy}和定理~\ref{thm:normbarrier}即可。
\end{proof}
这里得到的是\emph{一个上界表达式的下界}，不是 $|\Cx|$ 的下界。
它不排除实际两系数之间的消去，也不排除其他预条件化或系数敏感方法。

\subsection{增长维数的负方向}
\begin{theorem}[低秩修补的限定障碍]\label{thm:rank}
置 $\tau_x=Q^2/\log Q$，
\[
J_x=\left\lfloor N\left(\frac1{4Q}-\frac1H\right)\right\rfloor,\quad
d_x=2J_x+1,\quad k_x=d_x-35.
\]
对全部充分大的 $x$，存在复子空间 $W_x$，使
\[
\dim_\CC W_x\ge k_x,\qquad
\ip f{Af}\le-\frac{\tau_x}{2}\norm f^2\quad(f\in W_x).
\]
特别地，任何复秩小于 $k_x$ 的矩阵 $E$，即使非 Hermitian，也满足
$\norm{A-E}\ge\tau_x/2$，且
$k_x=Q^2/4-x^{11/32}+O(1)$。
\end{theorem}
\begin{proof}
令 $U_0=A+\diag\Delta\ge0$，并取
$f_j(n)=N^{-1/2}\e(jn/N)$，$|j|\le J_x$。
这些向量在 $N$ 个连续整数上正交；最终 $d_x<N$。
对全部 $q\le2Q$ 与 $|v|\le1$，
$q|j/N+v/H|\le1/2$。前证的剩余类方差界于是给出
\[
\ip{f_j}{U_0f_j}\le\frac{C_0}{N}\sum_q q^2,\qquad C_0=(\pi+1)^2.
\]
令 $F$ 以这些向量为列，$B=F^*U_0F\ge0$，$M_1=\sum_q q$，则
\[
\tr B\le C_0(1+2Q/N)M_1.
\]
另一方面，任一实际整数至多被两个壳层素数整除，所以
\[
\Delta(n)\ge D_*:=\sum_q\frac{q(q-2)}{q-1}-4Q
\sim\frac32\tau_x .
\]
$B$ 中大于 $D_*/2$ 的特征值数目至多 $2\tr B/D_*<35$，最终同时有
$D_*\ge\tau_x$。取 $B$ 的 $[0,D_*/2]$ 谱子空间经 $F$ 映射后的空间为 $W_x$。
它丢失至多 $35$ 个维度，并且其 $U_0$ 二次型不超过 $D_*/2$，
而被删除对角至少为 $D_*$。这证明负性断言。
若 $\rank E<k_x$，则 $W_x\cap\ker E$ 非零；在其中取单位向量即得
$\norm{A-E}\ge|\ip f{(A-E)f}|\ge\tau_x/2$。
用 $N=x/2+O(1)$ 展开 $d_x$ 得到维数渐近式。
\end{proof}
这一谱截断不意味着整个低频带都为负，也不声称某个更高秩阈值以上就能成功。

若要求实见证，不能直接把复谱子空间改称实空间。
在同一对称频带改用 $f_0$、$(f_j+f_{-j})/\sqrt2$、
$(f_j-f_{-j})/(\ii\sqrt2)$ 组成实正交矩阵 $V$。
实对称矩阵 $\operatorname{Re}(V^*U_0V)$ 正半定，且迹等于 $\tr B$。
对它重复谱截断，得到不同的实子空间，维数仍至少 $k_x$。
具体地，对于实线性映射
\[
E:\RR^{I_x}\longrightarrow\CC^{I_x},\qquad \rank_{\RR}E<k_x,
\]
在该实负空间与 $\ker E$ 的交中取单位向量，便得到
$\|A-E\|_{\RR^{I_x}\to\CC^{I_x}}\ge\tau_x/2$。
这里范数使用实定义域的 Euclidean 单位球，秩指实像空间的维数，不能与复秩混用。
等式 $f^TAf=f^T\operatorname{Re}(A)f$ 只适用于同一个实向量的二次型，
不能用来把两个不同实际系数的 $w^TA\beta$ 替换成 $w^T\operatorname{Re}(A)\beta$。
